\chapter{Dal topo al primate: i nostri cugini cognitivi}

\epigraph{``Se pensate che i bias cognitivi siano una peculiarità umana, vi sbagliate di grosso. Anche i topi hanno i loro pregiudizi, e stranamente assomigliano molto ai nostri.''}{}

\section{Il topo paranoico batte il topo filosofo}

Negli esperimenti di laboratorio, i topi mostrano chiari segni di ``bias di conferma'': se hanno imparato che un certo suono precede il cibo, continueranno a cercarlo anche quando il suono non è più correlato al premio\footnote{Rescorla, R. A. (1988). Pavlovian conditioning: It's not what you think it is. \emph{American Psychologist}, 43(3), 151-160.}. Proprio come noi continuiamo a controllare compulsivamente lo smartphone anche quando sappiamo che probabilmente non ci sono nuovi messaggi interessanti.

I piccioni sono ancora più affascinanti. In esperimenti dove ricevono cibo casualmente, sviluppano ``rituali superstiziosi'': se per caso stavano girando a sinistra quando è arrivato il cibo, continueranno a girare a sinistra nella speranza che si ripeta\footnote{Skinner, B. F. (1948). 'Superstition' in the pigeon. \emph{Journal of Experimental Psychology}, 38(2), 168-172.}. Qualcuno ha mai sentito parlare di rituali portafortuna negli umani?

La verità è che questi ``errori'' cognitivi sono così vantaggiosi che li ritroviamo in quasi tutto il regno animale. Un topo che assume che ogni rumore sia pericoloso vivrà più a lungo di uno che vuole sempre verificare la fonte del rumore. Anche se questo significa scappare inutilmente il 90\% delle volte.

Ma fermiamoci un attimo su questa parola: ``errori''. È davvero un errore quello che fanno questi animali? O forse stiamo guardando la cosa dal punto di vista sbagliato?

Quello che emerge dagli studi sul comportamento animale è una verità scomoda per il nostro ego di specie ``superiore'': i bias cognitivi non sono il prezzo che paghiamo per avere un cervello complesso. Sono il motivo per cui abbiamo un cervello complesso. Sono la base stessa dell'intelligenza, non una sua imperfezione.

Pensiamo al topo che ha paura di ogni rumore. Dal nostro punto di vista umano, sembra paranoico e irrazionale. Ma dal punto di vista evolutivo, quel topo ha risolto brillantemente un problema di ottimizzazione: come massimizzare le probabilità di sopravvivenza in un mondo pieno di incertezze.

Il topo non può permettersi di fare un'analisi statistica approfondita ogni volta che sente un rumore. Non ha tempo di valutare probabilità, consultare database storici, o fare brainstorming con altri topi. Deve decidere in una frazione di secondo: scappare o non scappare? E la strategia che funziona meglio, statisticamente parlando, è: quando in dubbio, scappa.

È la stessa logica che abbiamo visto nei capitoli precedenti, ma qui la vediamo nella sua forma più pura, non complicata dalle sovrastrutture culturali e sociali umane. Il cervello animale è un organo per prendere decisioni rapide in condizioni di incertezza, non per raggiungere la verità filosofica.

\section{L'intelligenza delle formiche e la stupidità delle folle}

Una formica singola è, diciamolo pure, piuttosto stupida. Non sa dove sta andando, non ha un piano generale, e se la separate dalla colonia probabilmente morirà presto. Eppure, migliaia di formiche insieme costruiscono opere architettoniche di precisione ingegneristica, trovano i percorsi più efficienti per il cibo, e mantengono società complesse che funzionano da milioni di anni\footnote{Hölldobler, B. \& Wilson, E. O. (1990). \emph{The Ants}. Cambridge, MA: Harvard University Press.}.

Questo paradosso -- intelligenza collettiva da stupidità individuale -- ci insegna qualcosa di fondamentale sul nostro cervello sociale. Anche noi, come le formiche, abbiamo evoluto potenti istinti per seguire il gruppo. E proprio come le formiche, a volte questi istinti ci portano a risultati brillanti, altre volte a disastri collettivi.

Il ``bias del gregge'' che ci fa seguire mode assurde o bolle finanziarie non è un difetto del nostro cervello sociale: è una caratteristica. Per la maggior parte della nostra storia evolutiva, fare quello che facevano tutti era la strategia di sopravvivenza più sicura. Se tutta la tribù scappava in una direzione, probabilmente c'era un buon motivo, anche se non lo capivi.

Il problema sorge quando questo istinto si scontra con il mondo moderno, dove ``tutti lo stanno facendo'' può significare ``tutti stanno per perdere i propri risparmi in una bolla speculativa''.

Ma per capire davvero come funziona l'intelligenza collettiva -- e quando diventa stupidità collettiva -- dobbiamo partire dal basso, dalle regole semplici che governano il comportamento di gruppo.

Le formiche non hanno un presidente, non tengono riunioni, non fanno piani strategici. Eppure riescono a coordinare attività che coinvolgono milioni di individui con una precisione e un'efficienza che farebbero invidia a qualsiasi multinazionale. Come ci riescono?

Il segreto sta in quello che i matematici chiamano ``emergenza'': comportamenti complessi che nascono dall'interazione di regole semplici\footnote{Johnson, S. (2001). \emph{Emergence: The Connected Lives of Ants, Brains, Cities, and Software}. New York: Scribner.}. Ogni formica segue solo tre regole di base: segui la traccia di feromoni più forte, lascia una traccia quando trovi cibo, esplora casualmente quando non trovi nessuna traccia.

Tre regole stupide, seguite da creature stupide, che insieme producono risultati di genio collettivo. È come se migliaia di persone che sanno fare solo addizioni semplici riuscissero insieme a risolvere equazioni differenziali.

Ma ecco il punto cruciale: questo sistema funziona perfettamente finché le regole di base rimangono appropriate per l'ambiente. Quando l'ambiente cambia più velocemente delle regole, l'intelligenza collettiva può trasformarsi rapidamente in stupidità collettiva.

\section{Primati: i nostri cugini troppo simili}

Gli scimpanzé condividono con noi il 98.8\% del DNA, e si vede. Nei test di laboratorio, mostrano molti dei nostri stessi bias cognitivi: avversione alle perdite, effetto ancoraggio, bias di conferma\footnote{Chen, M. K., Lakshminarayanan, V., \& Santos, L. R. (2006). How basic are behavioral biases? Evidence from capuchin monkey trading behavior. \emph{Journal of Political Economy}, 114(3), 517-537.}. È come guardare noi stessi allo specchio, solo con più peli e meno pretese di razionalità.

Ma c'è una differenza cruciale. Uno scimpanzé che ha paura irrazionale di un certo tipo di frutto al massimo avrà una dieta un po' limitata. Un umano che ha paura irrazionale degli OGM può influenzare le politiche alimentari di intere nazioni.

Il nostro problema non è avere un cervello di scimmia. Il problema è avere un cervello di scimmia con il potere di costruire armi nucleari, manipolare l'economia globale, e influenzare milioni di persone attraverso i social media.

Abbiamo la neurobiologia di una specie che viveva in gruppi di 50 individui, ma il potere di una specie che domina un pianeta. È un po' come dare le chiavi di una Ferrari a qualcuno che ha imparato a guidare sul trattore del nonno.

Ma andiamo con ordine, e vediamo quanto siamo davvero simili ai nostri cugini pelosi. Gli esperimenti sui primati degli ultimi decenni hanno rivelato somiglianze così inquietanti che a volte ci chiediamo se non stiamo semplicemente studiando versioni più semplici di noi stessi.

Prendiamo l'effetto ancoraggio, quel bias che fa sì che il primo numero che sentiamo influenzi tutte le nostre valutazioni successive. Negli esperimenti, agli scimpanzé viene mostrato un numero casuale, poi viene chiesto loro di scegliere tra diverse quantità di cibo\footnote{Beran, M. J. (2008). Monkeys (\emph{Macaca mulatta} and \emph{Cebus apella}) track, enumerate, and compare multiple sets of moving items. \emph{Journal of Experimental Psychology: Animal Behavior Processes}, 34(1), 63-74.}. Incredibilmente, le loro scelte sono sistematicamente influenzate dal numero che hanno visto per primo, esattamente come succede agli esseri umani.

È come se il cervello dei primati -- il nostro incluso -- fosse programmato per usare qualsiasi informazione disponibile come punto di riferimento, anche quando quella informazione è completamente irrilevante. È un meccanismo che probabilmente serviva per fare stime rapide nell'ambiente ancestrale: se vedi tre frutti su un albero, ti aspetti di trovarne un numero simile su alberi dello stesso tipo.

\subsection{L'avversione alle perdite nei primati}

Ancora più impressionante è come i primati gestiscono le perdite. Negli esperimenti di avversione alle perdite, ai bonobo vengono offerte due scelte: una ricompensa sicura di due frutti, o una ricompensa incerta che può essere di quattro frutti o zero frutti. Il valore atteso della seconda opzione è migliore (due frutti in media), ma i bonobo preferiscono sistematicamente la certezza di due frutti.

È esattamente il comportamento umano: preferiamo la certezza di un piccolo guadagno all'incertezza di un guadagno potenzialmente maggiore. Ma quando lo stesso esperimento viene fatto in termini di perdite -- perdere sicuramente due frutti o rischiare di perderne quattro o zero -- improvvisamente i bonobo diventano giocatori d'azzardo, proprio come noi.

Questo non è apprendimento o cultura: è neurobiologia pura. È il modo in cui il cervello dei primati -- tutti i primati, noi inclusi -- valuta rischi e ricompense. È stato programmato così da milioni di anni di evoluzione, e funziona secondo la stessa logica sia nel cervello di un bonobo che in quello di un trader di Wall Street.

\subsection{I bias sociali dei primati}

Ma forse l'aspetto più affascinante -- e preoccupante -- della nostra somiglianza con gli altri primati riguarda i bias sociali. Gli scimpanzé mostrano chiari segni di pregiudizio verso il gruppo di appartenenza: preferiscono sistematicamente membri del loro gruppo, anche quando questi si comportano peggio di membri di altri gruppi\footnote{Mahajan, N., Martinez, M. A., Gutierrez, N. L., Diesendruck, G., Banaji, M. R., \& Santos, L. R. (2011). The evolution of intergroup bias: Perceptions and attitudes in rhesus macaques. \emph{Journal of Personality and Social Psychology}, 100(3), 387-405. \textbf{Nota: Questo articolo è stato ritrattato nel 2014 dagli autori per problemi metodologici.}}.

In esperimenti controllati, gli scimpanzé aiutano più volentieri membri del loro gruppo che estranei, condividono più cibo con chi gli assomiglia, e mostrano più aggressività verso chi è diverso. Hanno perfino qualcosa che assomiglia al razzismo: discriminano sistematicamente contro scimpanzé di altre popolazioni, anche quando non c'è alcuna competizione per le risorse.

È lo stesso pattern che ritroviamo negli umani, solo amplificato. Noi non discriminiamo solo contro altre popolazioni: discriminiamo contro altre razze, altre nazioni, altre religioni, altri partiti politici, perfino altre squadre di calcio. Il meccanismo neurobiologico è identico, ma le sue manifestazioni sono diventate infinitamente più complesse e articolate.

\subsection{Dal topo al primate: una continuità evolutiva}

Quello che emerge da questa analisi comparativa è una continuità sorprendente nei bias cognitivi attraverso il regno animale. Dal topo che sviluppa superstizioni, alle formiche che seguono ciecamente le tracce di feromoni, agli scimpanzé che mostrano avversione alle perdite, fino a noi umani che compiamo gli stessi errori su scala planetaria.

Non sono difetti casuali o isolati: sono strategie cognitive che si sono dimostrate vincenti per milioni di anni. Il problema è che ora le stiamo applicando in contesti per cui non sono state progettate.

Un topo paranoico sopravvive meglio di uno filosofo. Una formica che segue il gruppo trova più cibo di una individualista. Uno scimpanzé che preferisce il suo gruppo ha più alleati di uno cosmopolita. E noi? Noi siamo l'apice di questa piramide evolutiva di bias ``di successo''.

La differenza non è che abbiamo superato questi meccanismi primitivi. La differenza è che li abbiamo elaborati, sofisticati, amplificati fino a livelli che i nostri cugini animali non potrebbero mai immaginare. Siamo topi paranoici con armi nucleari, formiche del gregge con i social media, scimpanzé tribali con ideologie globali.

E forse, riconoscendo questa continuità evolutiva, possiamo sviluppare un po' più di umiltà riguardo alla nostra presunta superiorità razionale. Dopo tutto, condividiamo gli stessi bias cognitivi di base con creature che consideriamo ``inferiori''. La differenza è solo nel grado di complessità con cui manifestiamo la stessa fondamentale irrazionalità.

In fondo, siamo tutti animali che cercano di sopravvivere usando strategie cognitive imperfette ma funzionali. Solo che alcuni di noi hanno smartphone.